\section{System Architecture}

The implemented system follows the two-stage retrieval architecture described in the proposed framework. The objective is to efficiently identify regulatory entities corresponding to extracted regulatory references while maintaining sufficient semantic accuracy for compliance-oriented applications.

The implemented architecture consists of four primary components:

\begin{enumerate}
    \item Regulatory reference processing and representation
    \item Dense candidate retrieval using a bi-encoder
    \item Semantic candidate reranking using a cross-encoder
    \item Confidence-based decision support
\end{enumerate}

The bi-encoder provides an efficient retrieval mechanism by mapping regulatory references and candidate entities into a shared embedding space. Retrieved candidates are subsequently evaluated by the cross-encoder, which performs deeper interaction-based relevance estimation.

\subsection{Model Components}

The retrieval component is implemented using a transformer-based bi-encoder architecture. The model receives regulatory references and candidate entities as independent inputs and generates dense vector representations used for similarity-based retrieval.

The reranking component is implemented using a transformer-based cross-encoder architecture. Unlike the bi-encoder, the cross-encoder jointly processes the regulatory reference and candidate entity, allowing direct token-level interaction through self-attention.

\subsection{Experimental Environment}

The experiments were conducted using the implemented training and evaluation pipeline. The system was developed using the PyTorch deep learning framework together with transformer-based model implementations. The computational environment used throughout model development and evaluation is summarized in Table~\ref{tab:experimental_environment}.

The computational environment was kept consistent across experiments to ensure comparable evaluation results.

\begin{table}[ht]
\centering
\caption{Experimental environment used for model development and evaluation.}
\label{tab:experimental_environment}
\begin{tabular}{ll}
\hline
\textbf{Component} & \textbf{Specification} \\
\hline
GPU & NVIDIA RTX 4070 (12GB VRAM) \\
CPU & AMD Ryzen 7 5800X \\
Memory & 32GB DDR4 RAM \\
Operating System & CachyOS Linux \\
Frameworks & PyTorch, Hugging Face Transformers \\
\hline
\end{tabular}
\end{table}

The computational environment was kept consistent across experiments to ensure comparable evaluation results.

\subsection{Model Selection and Configuration}

The proposed framework consists of multiple transformer-based components, each serving a distinct role within the regulatory entity linking pipeline. Table~\ref{tab:model_configuration} summarizes the pretrained models used as initialization points and their respective purposes within the framework.

\begin{table}[H]
\centering
\caption{Models used within the implemented regulatory entity linking framework.}
\label{tab:model_configuration}
\begin{tabular}{ll}
\hline
\textbf{Component} & \textbf{Base Model} \\
\hline

Embedding Model &
google/embeddinggemma-300m \cite{embeddinggemma2025}
\\[0.2cm]

Cross-Encoder &
BAAI/bge-reranker-v2-m3 \cite{chen2024bge}
\\[0.2cm]

Large Language Model &
\makecell[l]{Gemma-4-12B-it \cite{gemma4technicalreport2026}\\
Unsloth UD-Q4\_K\_XL GGUF Quantization \cite{unslothgemma4gguf2026}}
\\

\hline
\end{tabular}
\end{table}

Figure~\ref{fig:model_architecture} provides an overview of the implemented system architecture and illustrates the interaction between the individual components of the regulatory entity linking pipeline.

The bi-encoder serves as the initial candidate retrieval component and is fine-tuned on regulatory reference-entity pairs to learn domain-specific semantic representations. The reranking model is subsequently fine-tuned to perform fine-grained relevance estimation between retrieved candidates and regulatory references. Additionally, a large language model is utilized for domain adaptation through noise characterization and data augmentation.

The system receives regulatory documents as input and extracts relevant regulatory references. These references are encoded into dense representations and compared against indexed regulatory entities to retrieve a set of candidate matches. The candidate set is subsequently evaluated by the cross-encoder, which produces relevance scores used for final entity linking decisions.

\begin{figure}[ht]
\centering

\begin{tikzpicture}[
    node distance=1.4cm,
    every node/.style={font=\small},
    block/.style={
        rectangle,
        rounded corners,
        draw=black,
        minimum width=5cm,
        minimum height=0.9cm,
        align=center,
        fill=gray!10
    },
    decision/.style={
        diamond,
        draw=black,
        aspect=2,
        align=center,
        fill=gray!10
    },
    arrow/.style={->, thick}
]

\node[block] (input)
{Regulatory Documents};

\node[block, below=of input] (extract)
{Regulatory Reference Extraction};

\node[block, below=of extract] (embedding)
{Bi-Encoder Retrieval\\
EmbeddingGemma-300M};

\node[block, below=of embedding] (candidate)
{Top-$k$ Candidate\\
Regulatory Entities};

\node[block, below=of candidate] (reranker)
{Cross-Encoder Reranking\\
BGE-reranker-v2-m3};

\node[block, below=of reranker] (score)
{Relevance Score\\
Confidence Estimation};

\node[decision, below=of score, yshift=-0.3cm] (threshold)
{Confidence\\
Threshold};

\node[block, below left=1.5cm and 1cm of threshold] (auto)
{Automatic\\
Entity Linking};

\node[block, below right=1.5cm and 1cm of threshold] (human)
{Human Expert\\
Review};

\draw[arrow] (input)--(extract);
\draw[arrow] (extract)--(embedding);
\draw[arrow] (embedding)--(candidate);
\draw[arrow] (candidate)--(reranker);
\draw[arrow] (reranker)--(score);
\draw[arrow] (score)--(threshold);

\draw[arrow] (threshold)--node[above left]{High confidence}(auto);
\draw[arrow] (threshold)--node[above right]{Low confidence}(human);

\end{tikzpicture}

\caption{Implemented regulatory entity linking architecture.}
\label{fig:implemented_architecture}

\end{figure}